\section{Ideal classes that survive a quadratic tower}
We construct the arithmetic curve by separating a finite-group
calculation from its realization by ray class fields.

\subsection{The group calculation}
Let $V$ be a nonzero finite-dimensional $\F_2$-vector space and put
$G=(V\oplus V)\rtimes\langle\tau\rangle$, where $\tau$ exchanges the
two summands. Its elements are denoted $(x,y;\epsilon)$, with
$\epsilon\in\F_2$. Write $H=V\oplus V$ and
$\Sigma(x,y;\epsilon)=x+y$.

\begin{lemma}\label{lem:wreath}
The kernel $A_0$ of $\Sigma$ is the elementary abelian group
$\F_2\oplus V$, with $(b,u)$ corresponding to $(u,u;b)$.
The section $v\mapsto(v,0;0)$ makes its $V$-action
\[
 v\cdot(b,u)=(b,u+bv).
\]
The map of regular modules
\[
 \phi:\F_2[V]\longrightarrow A_0,\qquad e_g\longmapsto(1,g)
\]
is surjective and $V$-equivariant. It induces a surjection
$\F_2[V]\rtimes V\twoheadrightarrow G$ over $V$. Moreover,
\[
 \Ver_{G,H}(x,y;\epsilon)=(x+y,x+y).
\]
\end{lemma}
\begin{proof}
The diagonal pairs are fixed by interchange, so multiplication on
$A_0$ is addition. Conjugating $(u,u;b)$ by $(v,0;0)$ gives
$(u+bv,u+bv;b)$. The translation $e_g\mapsto e_{g+v}$ therefore
commutes with $\phi$. Its image contains $(1,0)$ and every
$(0,g)=(1,g)+(1,0)$, proving surjectivity. The images $(g,g;1)$
are conjugates of $\tau$; each has order two and its subgroup meets
$H$ only in the identity.

For transfer choose representatives $1,\tau$ for $G/H$. If
$\epsilon=0$, transfer multiplies $(x,y;0)$ by its conjugate and
gives $(x+y,x+y)$. If $\epsilon=1$, the two cosets are exchanged,
so transfer is the square, with the same value. In the alternative
coordinates $(b,u;v)\in A_0\rtimes V$, this also gives
\begin{equation}\label{eq:square}
 (b,u;v)^2=(0,bv;0).
\end{equation}
\end{proof}

For $V=\F_2$, $G$ has order eight. The two regular basis vectors
map to the two reflections $(0,0;1)$ and $(1,1;1)$; their sum is
the central diagonal element. Transfer sends either off-diagonal
element of $H$ to that central element. This is the smallest case
of the class-preserving map in the next lemma.

\subsection{One arithmetic step}
For a totally imaginary number field $F$, set
$C_F=\Pic(\OO_F[1/2])/2$. This is a finite-dimensional
$\F_2$-vector space.

\begin{lemma}\label{lem:step}
Suppose $r=\dim C_F>0$, and let $S$ be a finite set of finite
primes of $F$ containing all dyadic primes. There is a quadratic
extension $L/F$ such that every prime in $S$ is unramified and
inert, extension of ideals is injective on $C_F$, and
$\dim C_L\ge2r$.
\end{lemma}
\begin{proof}
Class field theory identifies $C_F$ with the Galois group $V$ of
the maximal elementary abelian two-extension $M/F$ unramified at
finite primes and split at dyadic primes
\cite[Chapter V, Theorems 3.5--3.6, p.~158]{MilneCFT}. Fix, over
each $\mathfrak q\in S$, one prime $\mathfrak Q$ of $M$. Prescribe
at that prime the unit squareclass giving the unramified quadratic
extension, and prescribe the squareclass of $1$ at all other
primes above $\mathfrak q$. Thus the nonsquare prescription has exactly
one chosen prime in each fiber.

The unramified quadratic class is a unit even at a dyadic prime:
an irreducible residual polynomial $T^2-T+\bar c$ lifts with
unit discriminant $1-4c$. Local squareclasses are open, by
Hensel's lemma. Weak approximation gives an element $b\in M^\times$
with all these local classes and with valuation zero at those
primes. Choose a $V$-invariant ray modulus $\mathfrak m$, supported
there and sufficiently deep to preserve the squareclasses.

Let $E/M$ be the ray class field. Since $V$ preserves the ray
subgroup, $E/F$ is Galois. In $\Gal(E/M)$ take the Artin image
of the ray class of $(b)$. Chebotarev gives a prime
$\mathfrak p\notin S$ of $F$ and a prime of $E$ whose Frobenius
is that element \cite[Chapter V, Theorem 3.23, p.~164]{MilneCFT}.
The Frobenius is trivial on $M$, so $\mathfrak p$ splits completely
in $M$. The prime $\mathfrak P$ of $M$ so obtained has ray class
$[(b)]$. Thus
\[
 \mathfrak P=(\eta),\qquad \eta/b\equiv1\pmod{\mathfrak m}
\]
for some integral $\eta\in M$. Indeed equality of ray classes
makes $\mathfrak P(b)^{-1}$ generated by an element congruent to
one, and equality of the resulting principal ideal with
$\mathfrak P$ makes its generator integral.

Adjoin all conjugate square roots:
\[
 \widetilde N=M(\sqrt{\sigma\eta}:\sigma\in V).
\]
Their squareclasses are independent: the valuation at
$\sigma\mathfrak P$ is odd only for $\sigma\eta$. Hence
$\Gal(\widetilde N/M)=\F_2[V]$. The automorphism $g\in V$
lifts by permuting the labelled generators according to
$\sqrt{\sigma\eta}\mapsto\sqrt{g\sigma\eta}$. Independence
ensures that the algebra with these generators is a field; the
permutations respect its defining relations and compose according
to the multiplication in $V$. Therefore
$\Gal(\widetilde N/F)=\F_2[V]\rtimes V$.

Quotient by $\ker\phi$ from Lemma~\ref{lem:wreath}. The resulting
field $N$ has group $G$, and its quotient $\Sigma$ is the original
extension $M/F$. Put $L=N^H$. This is quadratic over $F$.

We check all possible ramification. Outside dyadic primes and the
primes over $\mathfrak p$, the radicands are units at odd residue
characteristic, so their quadratic extensions are split or
unramified. At dyadic primes this follows from the prescription.
At $\sigma\mathfrak P$ the inertia in $\widetilde N/M$ is generated
by $e_\sigma$. Its image in $G$ is a conjugate of $\tau$ and has
trivial intersection with $H$. Since $\mathfrak p$ splits in
$M/F$, the same computation shows $N/L$ is unramified there.
Thus $N/L$ is unramified at every finite prime. There are no real
places.

For the inertness assertion, let $v\in V$ be the Frobenius at
$\mathfrak q\in S$. If $v=0$, the local Frobenius in the regular
module is one $e_g$, because exactly one conjugate radicand has
the nonsquare unramified class. Its image $(1,g)$ has order two
and maps nontrivially to $G/H$. Hence $\mathfrak q$ is inert in
$L/F$, and the prime above it splits in $N/L$. This applies in
particular at every dyadic prime.

If $v\ne0$, its decomposition group in $M/F$ is $\langle v\rangle$.
At a chosen completion exactly two conjugate radicands have the
nonsquare class, indexed by $g$ and $g+v$. The local Frobenius over
$M$ consequently maps to $\phi(e_g+e_{g+v})=(0,v)$. If $z$ is
Frobenius over $F$, then $\Sigma(z)=v$ and $z^2=(0,v)$.
Equation~\eqref{eq:square} forces the first coordinate $b$ of
$z=(b,u;v)$ to be one. It therefore maps nontrivially to $G/H$,
again proving unramified inertness.

The unramified, dyadically split elementary abelian extension
$N/L$ gives a surjection $C_L\twoheadrightarrow V\oplus V$.
Compatibility of reciprocity with extension of ideals is transfer
\cite[Chapter V, diagram and paragraph after Proposition 3.18,
p.~162]{MilneCFT}. The composite
of $C_F\to C_L$ with that surjection is therefore
$v\mapsto(v,v)$ under $C_F=V$. It is injective, proving the
injectivity of $C_F\to C_L$. The quotient of dimension $2r$
proves the required dimension bound.
\end{proof}

\subsection{The limit ring and its cohomology}
We now prove Theorem~\ref{thm:arithmetic}, including the properties
that can fail for an uncontrolled infinite algebraic tower.

\begin{proof}[Proof of Theorem~\ref{thm:arithmetic}]
Start with $F_0=\Q(\sqrt{-35})$ and
$w=(1+\sqrt{-35})/2$, so $w^2-w+9=0$ and the discriminant is
$-35$. Minkowski's bound is $(2/\pi)\sqrt{35}<4$
\cite[Theorem 4.3, pp.~70--71]{MilneANT}. The prime 2 is inert
and principal. The two primes of norm 3 are
$P=(3,w)$ and $\bar P=(3,w-1)$. Since $N(w)=9$ and $w$ is
a unit at $\bar P$, one has $(w)=P^2$. There is no element of
norm 3: for integers $x,y$,
\[
 N(x+yw)=x^2+xy+9y^2
       =(x+y/2)^2+35y^2/4.
\]
For $y\ne0$ this exceeds 3, and for $y=0$ it is a square.
Minkowski's bound now proves $\Pic(\OO_{F_0})=\Z/2$, with
generator $P$. Inverting the principal dyadic prime leaves it
unchanged, so $\dim C_{F_0}=1$.

Enumerate the rational primes with 2 first. At stage $n$ prescribe
inertness at every prime of $F_n$ above the first $n+1$ primes,
and use Lemma~\ref{lem:step} to choose $F_{n+1}$. Define
\[
 K=\bigcup_nF_n,\qquad R=\bigcup_n\OO_{F_n}[1/2].
\]
This is the integral closure of $\Z[1/2]$ in $K$. For any fixed
rational prime, splitting stops after a finite stage and all
subsequent steps are unramified. There are consequently finitely
many primes of $R$ over it, with bounded ramification indices.
For a nonzero prime $P$ of $R$, the local ring $R_P$ is the
union of the compatible number-field valuation rings. Bounded
ramification makes its value group discrete, so it is a DVR.
Each nonzero element of $R$ is contained in only finitely many
maximal ideals: its number-field norm has finitely many rational
prime divisors and each splits only finitely often.

These facts imply noetherianity directly. For an ideal $I\ne0$,
choose $a\ne0$ in $I$, and let $P_1,\ldots,P_s$ be the maximal
ideals containing $a$. At each $P_j$ choose $b_j\in I$ with
the minimum possible valuation. The ideal
$J=(a,b_1,\ldots,b_s)$ has the same localization as $I$ at every
maximal ideal, and hence $I=J$. Thus $R$ is a Dedekind domain.
This is the finite-splitting criterion of Popescu--Vraciu
\cite[pp.~40--42]{PopescuVraciu}, here proved in the form needed.
A Dedekind domain with characteristic-zero fraction field is
excellent \cite[Tag 07QS, Proposition 15.53.3(4)]{Stacks}.
The eventual-inertness construction of a Dedekind integral closure
also appears in \cite[\S5, pp.~511--512]{Popescu2022}; the
class-preserving extension lemma is the additional requirement here.

Let $T/K$ be finite. It is $EK$ for a number field $E$ containing
some $F_n$. Any element of $\Br(T)[2]$ descends, after increasing
$E$, to an element of $\Br(E)[2]$: finite presentation descends
the algebra, and a later finite stage descends the relation that
its double is zero. This class has finitely many nonzero local
invariants. Increase $n$ so that all underlying rational primes
are prescribed, replacing $E$ by $EF_n$ and the class by its
restriction if necessary. At each such prime, the local degrees in
$F_m/F_n$ are $2^{m-n}$. In a common completion, for any prime
$w_m$ of $EF_m$ over $w$ of $E$, the degree formula gives
\[
 v_2([(EF_m)_{w_m}:E_w])
 \ge m-n-v_2([E_w:(F_n)_v]).
\]
Thus, for one sufficiently large $m$, all the relevant degrees
are even, uniformly over their prolongations. Restriction multiplies
local Brauer invariants by these degrees. All archimedean places
are complex. Brauer--Hasse--Noether therefore kills the class over
$EF_m$ \cite[Chapter III, Theorem 2.1, p.~103, and Chapter VIII,
Theorem 4.2, p.~241]{MilneCFT}.
We have shown $\Br(T)[2]=0$ for every finite extension $T/K$.

Choose a maximal pro-two subgroup of the absolute Galois group of
$K$, and let $K'/K$ be its fixed field. Kummer theory and continuity
over the finite subextensions give
\[
 H^2(K',\F_2)=H^2(K',\mu_2)
 =\varinjlim_{K\subset T\subset K',\,[T:K]<\infty}\Br(T)[2]=0.
\]
The pro-prime subgroup criterion for Galois cohomological dimension
\cite[Tag 0DV9, Lemma 59.67.3]{Stacks} now gives vanishing in
every degree at least two for all two-primary torsion coefficients.
Thus $\cd_2K\le1$. There is a nonzero prime of $R$
by lying over an odd prime of $\Z[1/2]$. A uniformizer of its
DVR has odd valuation and is not a square in $K$, so Kummer
theory proves $\cd_2K=1$.

Invertible modules and their isomorphisms descend through filtered
ring colimits. Hence
\[
 \Pic(R)/2=\varinjlim_n C_{F_n}.
\]
Every transition is injective and $\dim C_{F_n}\ge2^n$.
This vector space is therefore infinite-dimensional. A nonempty
open of the noetherian curve $\Spec R$ deletes finitely many
closed points. Its Picard group is the quotient by their divisor
classes, so its reduction modulo two still has infinite dimension.
Kummer theory gives $H^2(U,\mu_2)\ne0$ for every such open.

Every closed residue field is an increasing union of finite fields
with successive degrees eventually two. Its absolute Galois group
has trivial pro-two subgroup, and its two-cohomological dimension
is zero. The generic field contains a totally imaginary number
field; thus no residue field is orderable. The generic-field bound
\[
 \cd_2(U)\le\cd_2(K)+\dim U=2
\]
follows from SGA~4, X, Corollary 4.4, p.~57 \cite{SGA4}.
Its strict-local hypothesis holds here: at a closed point the
strict local ring is a DVR with separably closed residue field
of odd characteristic, so its fraction field has two-cohomological
dimension one by X, Theorem 2.2, p.~46; for a zero-dimensional
irreducible closed subscheme the strict localization has
separably closed residue field and dimension zero. Thus the bound
applies to arbitrary two-primary torsion sheaves. Together with
the Kummer lower bound this proves $\cd_2(U)=2$. At a closed
point the residue dimension plus local-ring dimension is $0+1$;
at the generic point it is $1+0$. This proves the final supremum.
\end{proof}
